\documentclass[UTF8,12pt]{ctexart}
\usepackage[a4paper,margin=2.4cm]{geometry}
\usepackage{graphicx}
\usepackage{booktabs}
\usepackage{longtable}
\usepackage{array}
\usepackage{enumitem}
\usepackage{xcolor}
\usepackage{float}
\usepackage{caption}
\usepackage{hyperref}
\usepackage{tikz}
\usetikzlibrary{arrows.meta,positioning,shapes.geometric,fit,calc}

\hypersetup{
  colorlinks=true,
  linkcolor=blue!55!black,
  urlcolor=blue!55!black
}

\graphicspath{{../screenshots/}}
\setlength{\parindent}{2em}
\setlength{\parskip}{0.35em}
\setlist{nosep,leftmargin=2em}
\renewcommand{\arraystretch}{1.25}
\captionsetup{font=small,labelfont=bf}
\sloppy

\definecolor{mainblue}{HTML}{1F5F99}
\definecolor{lightblue}{HTML}{EAF2FF}
\definecolor{linegray}{HTML}{CED6E0}
\definecolor{softyellow}{HTML}{FFF6DF}

\newcommand{\code}[1]{\texttt{\detokenize{#1}}}
\newcommand{\docmeta}[2]{\noindent\textbf{#1：}#2\par}
\newcommand{\screenshot}[3][0.8\textwidth]{
  \begin{figure}[H]
    \centering
    \includegraphics[width=#1]{#2}
    \caption{#3}
  \end{figure}
}

\tikzset{
  flow/.style={rectangle,rounded corners,draw=mainblue,fill=lightblue,minimum width=2.6cm,minimum height=0.9cm,align=center,font=\small},
  decision/.style={diamond,aspect=2,draw=mainblue,fill=softyellow,align=center,font=\small,inner sep=1pt},
  dataobj/.style={rectangle,draw=linegray,fill=white,minimum width=2.6cm,minimum height=0.85cm,align=center,font=\small},
  actor/.style={rectangle,rounded corners,draw=mainblue,fill=lightblue,minimum width=2.2cm,minimum height=0.9cm,align=center,font=\small\bfseries},
  usecase/.style={ellipse,draw=mainblue,fill=white,minimum width=2.9cm,minimum height=0.85cm,align=center,font=\small},
  classbox/.style={rectangle,draw=mainblue,fill=white,minimum width=3.1cm,minimum height=1.2cm,align=left,font=\small},
  arrow/.style={-{Latex[length=3mm]},thick,draw=mainblue},
  line/.style={thick,draw=mainblue}
}



\title{\textbf{作业管理系统过程剖析}}
\author{}
\date{}

\begin{document}
\maketitle

\section{问题一：初始化脚本导入失败}

\subsection{问题现象}

在项目根目录执行数据库初始化脚本时，终端出现 \code{ModuleNotFoundError}，提示找不到 \code{assignment_system} 模块。

\screenshot{bug_01_import_error_flat.jpg}{初始化脚本导入失败}

\subsection{排查过程}

开始以为是虚拟环境没有安装依赖，但检查后发现 Flask 已经安装。继续观察错误位置，发现报错发生在 \code{scripts/init_db.py} 第一行导入项目包时。原因是直接运行 \code{scripts/init_db.py} 时，Python 会把 \code{scripts} 目录放到模块搜索路径最前面，而不是项目根目录，所以找不到同级目录外的 \code{assignment_system} 包。

\subsection{修复方法}

在脚本开头计算项目根目录，并将根目录加入 \code{sys.path}，然后再导入系统模块。修复后再次执行初始化脚本，数据库可以正常创建。

\newpage

\section{问题二：启动脚本重复重建数据库}

\subsection{问题现象}

检查 Windows 一键运行流程时发现，\code{run_windows.bat} 每次启动都会执行 \code{python scripts/init_db.py}。原来的初始化脚本会直接调用 \code{init_db(seed=True)}，而 \code{init_db} 会执行 \code{schema.sql} 中的 \code{DROP TABLE} 语句。这样第二次启动系统时，之前提交的作业和评分记录都会被清空。

\screenshot{bug_02_init_reset_flat.jpg}{重复初始化风险}

\subsection{排查过程}

先查看 \code{run_windows.bat}，确认启动前一定会调用初始化脚本。再查看 \code{scripts/init_db.py}，发现脚本没有区分“首次建库”和“重置数据库”。最后查看 \code{schema.sql}，发现文件开头包含多个 \code{DROP TABLE IF EXISTS} 语句，所以问题根源是启动脚本把“重置数据库”当成了“检查数据库”。

\subsection{修复方法}

新增 \code{ensure_db} 方法，只检查数据库中是否存在 \code{users} 表。如果表不存在，说明是第一次运行，再创建数据库和演示数据；如果表已经存在，就保留原有数据。初始化脚本增加 \code{--reset} 参数，只有手动执行 \code{python scripts/init_db.py --reset} 才会清空并重建数据库。

\newpage

\section{功能验证过程}

修复后进行了以下验证：

\begin{enumerate}
  \item 管理员账号可以登录后台。
  \item 待审核教师显示在用户表中，管理员可以审核。
  \item 教师可以创建课程并发布作业。
  \item 学生可以加入课程并提交作业。
  \item 教师可以查看提交记录，填写成绩和评语。
  \item 教师可以导出 CSV 统计文件。
  \item 自动化测试通过。
\end{enumerate}

测试命令：

\begin{verbatim}
.venv/bin/python -m unittest discover -s tests -v
\end{verbatim}

测试结果：

\begin{verbatim}
Ran 2 tests in 1.293s
OK
\end{verbatim}

\screenshot{08_teacher_after_grade.png}{评分后页面}

\newpage

\section{总结}

本次开发过程中，主要问题集中在运行环境和数据持久化两个方面。第一个问题说明脚本直接运行时要注意模块搜索路径；第二个问题说明初始化逻辑不能随意清空数据库，尤其是答辩现场需要多次启动系统时，更要保证数据不会丢失。

通过这两个问题的排查，系统最终形成了更稳定的运行方式：源码运行时自动检查数据库，必要时才初始化；如需重新演示初始数据，可以手动执行重置参数。这样既能保证首次运行方便，也能避免答辩时误删数据。

\end{document}
